\documentclass[10pt]{article}

\usepackage[utf8]{inputenc}

\usepackage[T1]{fontenc}

\usepackage{amsmath}

\usepackage{amsfonts}

\usepackage{amssymb}

\usepackage[version=4]{mhchem}

\usepackage{stmaryrd}

\usepackage{graphicx}

\usepackage[export]{adjustbox}

\graphicspath{ {./images/} }

\usepackage{bbold}

\usepackage{mathrsfs}



\begin{document}

лет, т. 1, М., 1959 , с. $383 \_98 ;$ [18] Г о н ч а р А. А., М е р г ел я н С. Н., в кн.: История отечественной математики, т.

кн. 1, К. 1970, с. $112-93 ;$ [19] Т а м р а $о$ в I. Гладкости и полиномиальные приближения, К., $1975 ;[20]$ М е льн и к о в М. С., С и нан я н С. О., внн.: Современные проблемы математики, т. 4, М., 1975 , с. $143-250$



ПРИБЛИЖЕНИЯ ПОРЯДОК, а п п р о к с и м ац и и п о р я д о к,- порядок погрешности приближения как переменной величины, зависящей от непрерывного или дискретного аргумента $\tau$, относительно другой переменной $\varphi(\tau)$, поведение к-рой, как правило, считается известным. Обычно $\tau$ - нек-рый параметр, являющийся числовой характеристикой приближающего множества, (напр., размерность этого множества) или метода приближения (напр., шаг интерполяции); при әтом множество значений $\tau$ имеет конечную или бесконечную предельную точку. Функция $\varphi(\tau)$ - чапе всего степенная, показательная или логарифмическая. В качестве $\varphi(\tau)$ может фигурировать непрерывности модуль приближаемой функции (или нек-рой ее производной) или его мажоранта.



П. п. характеризует как аппроксимативные возможности метода приближения, так и определенные свойства приближаемого обтекта, напр., дифференциальноразностные свойства приближаемой функции (см. $\Pi p u-$ ближение функций; прямые и обратные теоремы).



В численном анализе П. п. численного метода, имеющего погрешность $O\left(h^{m}\right)(h-$ шаг метода), наз. показатель $m$.



Лum.: [1] Г о н ч а р о в В. Л., Теория интерполирования и приближения функций, 2 изд., М., 1954; [2] Т и м а н А. Ф., Теория приближения функций пействительного переменного,



\begin{center}

\includegraphics[max width=\textwidth]{2023_07_08_35d33b1f96a32a44582fg-1}

\end{center}



ПРИБЛИЖЕНИЯ ТЕОРИЯ, $\quad$ п п р о к с п м а ц и и т е о р и я,- раздел математич. анализа, изучающий методы приближения одних математич. об'ьктов другими и вопросы, связанные с исследованием и оценкой возникающей при этом погрешности.



Основное содержание П. т. относится к приближению функций. Фундамент П. т. был заложен работами П. Л. Чебышева (1854-59) о наилучшем равномерном приближении функций многочленами и К. Вейерптрасca (K. Weierstraß), установившего в 1885 приндипиальную возможность приблизить непрерывную на конечном отрезке функцию алгебраич. многочленами со сколь угодно малой наперед заданной погрешностью. Развитие П. т. в значительной степени определялось основополагающими работами А. Jебега (H. Lebesgue), II. Ж. Валле Пуссена (Ch. J. La Vallée Poussin), C. Н. Бернштейна, Д. Джексона (D. Jackson), Ж. Фавара (J. Favard), А. Н. Колмогорова, С. М. Никольского о приближении функций и классов функций.



С развитием функционального анализа многие вопросы П. т. стали рассматривать в самой обцей ситуации, напр. как приближение элементов произвольного линейного нормированного пространства $X$. При этом выделяют три круга задач, к-рые в определенной степени соответствуют и основным хронологич. этапам развития исследований в П. т.



\begin{enumerate}

  \item Приближение фиксированного әлемента $x \in X$ әлементами фиксированного множества $\mathfrak{\Re} \subset X$. Если в качестве меры приближения взять величину

\end{enumerate}



\[

E(x, \mathfrak{l})=\inf _{u \in \mathfrak{R}} \| x-u \rrbracket

\]



т. е. наилучшее приближение $x$ множеством $\mathfrak{R}$, то, помимо исследования и оценки $E(x, \mathfrak{R})$, возникают вопросы о существовании элемента наилучішего приближения $u_{0}$ из $\mathfrak{A}$ (для к-рого $\left.\left\|x-u_{0}\right\|=E(x, \mathfrak{l})\right)$, его единственности и характеристич. свойствах. Јюбой оператор $A$, отображаюший $X$ в $\Re$, задает нек-рый метод приблияения с погрешностыо $\|x-A x\|$. Если $\mathfrak{N}-$ линейное многообразие, то особое значение имеют линей- ные операторы. Для последовательности $\left\{A_{n}\right\}$ таких операторов возникает задача об условиях сходимости $A_{n} x \rightarrow x$ для любого $x \in X$.



\begin{enumerate}

  \setcounter{enumi}{1}

  \item Приближение фиксированного множества $\mathbb{D} \subset X$ элементами другого фиксированного множества $\mathfrak{R}$ из $X$. Наилучшее приближение в әтом случае выражается величиной

\end{enumerate}



\[

E(\mathfrak{M}, \mathfrak{N})=\sup _{x \in \mathfrak{M}} E(x, \mathfrak{l})

\]



к-рая дает минимально возможную оценку погрешности приближения любого элемента $x \in X$ элементами из $\mathfrak{A}$. В конкретных случаях задача состоит в том, чтобы оценить или точно выразить $E(\mathfrak{N}), \Re)$ через характеристики, задающие множества ฏ) и ‥ Если приближение осуществляется с помощыо оператора $A$, то исследуется верхняя грань



\[

\sup _{x \in \mathfrak{D i}_{i}}\|x-A x\|

\]



а также (если $\mathfrak{N}$ - линейное многообразие) величина



\[

\mathscr{O}(\mathfrak{M}, \mathfrak{\Re})=\inf _{A X \subset \mathfrak{R}} \sup _{x \in \mathfrak{M}}\|x-A x\|,

\]



где нижняя грань распространена на все линейные операторы, отображающие $X$ в $\mathfrak{N}$. Линейный оператор, реализующий эту нижнюғ грань (если он существует), определяет наилучший линейный метод приближения. Особый интерес представляет выяснение случаев, когда $\mathfrak{O}(\mathfrak{M}, \mathfrak{R})=E(\mathfrak{M}, \mathfrak{N})$.



\begin{enumerate}

  \setcounter{enumi}{2}

  \item Наилучшее приближение фиксированного множества $\mathfrak{M} \subset X$ задавным классом $\{\mathfrak{Y}\{\}$ аппроксимирующих множеств из $X$. Предполагается, что в класо $\{\mathfrak{P}\}$ входят в каком-то смысле "равноценные» множества, напр. содержащие одно и то же количество элементов или имөющие одну и ту же размерность. Первый случай приводит к задаче об є-энтропии множества $\mathfrak{M}$ (относительно $X$ ), второй - к задачам вычисления noneречников множества $\mathfrak{M}$ (в пространстве $X$ ), в частности величин

\end{enumerate}



и



\[

d_{N}(\mathfrak{M}, X)=\inf _{\mathfrak{R}_{N}} E\left(\mathfrak{M}, \mathfrak{\Re}_{N}\right)

\]



\[

d^{\prime}{ }_{N}(\mathfrak{M}, \quad X)=\inf _{\mathfrak{R}} \mathscr{C}\left(\mathfrak{M}, \mathfrak{R}_{N}\right),

\]



где нижние грани берутся по всем подпространствам $\mathfrak{\Re}_{N}$ из $X$ фиксированной размерности $N$ (или по всевозможным их сдвигам $\left.\mathfrak{l}_{N}+a\right)$. Таким образом, в задачах (1)-(2) речь идет об отыскании наилучшего (соответственно наилучшего линейного) аппарата приближения размерности $N$ для множества $\mathfrak{M}$.



Лum [1] А х и е з е р Н. И., Лекции по теории аппроксимации, 2 изд , М., 1965 , [2] терполирования и приближения функций, 2 изд., М., 1954; [3] Д з я д ы к В. К., Введение в теорию равномерного приближения функций полиномами, М, 1977, [4] Ни к о ль ск и й С. М., Приближение функций многих переменных и теоремы вложения, 2 изд., М., 1977; [5] К о р н е й ч у к Н. П., Экстремальные задачи теории приближения, М., 1976; [6] Т и х о м иp o в B. М., Н екоторые вопросы теории приближений, M., вительного переменного, М., 1960.



Н. П. Корнейчук, В. П. Моторный ПРИБЛИЖЕНИЯ ФУНКЦИИ МЕРА - количественное выражение погрешности приближения. Когда речь идет о приближенци функции $f(t)$ функцией $\varphi(t)$ мера приближения $\mu(f, \varphi)$ обычно определяется метрикой нек-рого функционального пространства, содержащего как $f(t)$, так и $\varphi(t)$. Напр., если функции $f(t)$ и $\varphi(t)$ непрерывны на отрезке $[a, b]$, часто пользуются равномерной метрикой пространства $C[a, b]$, т. е. полагают



\[

\mu(f, \varphi)=\max _{a \leqslant t \leqslant b}|f(t)-\varphi(t)| .

\]





\end{document}